% !TeX root = main.tex
\chapter{Estructura del Proyecto}
\label{ch:estructura}

Este capítulo describe la organización de directorios del proyecto SIPREH en el frontend y el backend.

\section{Estructura del frontend}
\label{sec:estructura_frontend}

El frontend es una aplicación Next.js 14 que sigue las convenciones del \textit{App Router}:

\begin{verbatim}
drought-frontend/
+-- public/
|   +-- watersheds/          # GeoJSON de cuencas hidrografica
|   +-- icons/               # Iconos de marcadores
+-- src/
|   +-- app/
|   |   +-- layout.tsx       # Layout raiz
|   |   +-- page.js          # Pagina principal (mapa + sidebar)
|   |   +-- login/page.tsx   # Pagina de autenticacion
|   |   +-- admin/page.tsx   # Panel de administracion
|   +-- components/
|   |   +-- map/             # Componentes del mapa Leaflet
|   |   +-- charts/          # Graficas uPlot y Canvas 2D
|   |   +-- sidebar/         # Paneles de historico y prediccion
|   |   +-- admin/
|   |       +-- dashboard/
|   |           +-- FilesSection.js   # Carga de archivos Parquet
|   |           +-- DatasetConfig.js  # Configuracion de datasets
|   |           +-- UserManager.js    # Gestion de usuarios
|   |           +-- SyncStatus.js     # Estado de sincronizacion R2
|   +-- utils/
|   |   +-- exporters.js     # Exportacion PNG y JSON
|   |   +-- lttb.js          # Algoritmo de diezmado LTTB
|   |   +-- colorScale.js    # Paleta de colores de severidad
|   +-- lib/
|       +-- api.js           # Cliente HTTP hacia el backend
|       +-- auth.js          # Gestion de tokens JWT
+-- .env.local               # NEXT_PUBLIC_API_URL
+-- next.config.js
+-- package.json
\end{verbatim}

\begin{remark}
La función de exportación (\texttt{src/utils/exporters.js}) genera imágenes PNG de 1800$\times$1100 px con encabezado institucional y metadatos, y archivos JSON con los datos de la consulta activa (columnas, fuente, resolución, sistema de referencia WGS84).
\end{remark}

\section{Estructura del backend}
\label{sec:estructura_backend}

El backend es una aplicación FastAPI organizada en capas funcionales:

\begin{verbatim}
drought-backend/
+-- app/
|   +-- main.py              # Punto de entrada FastAPI
|   +-- config.py            # Variables de entorno y settings
|   +-- database.py          # Inicializacion SQLAlchemy
|   +-- models/              # Modelos ORM (usuario, dataset, archivo)
|   +-- api/
|   |   +-- v1/
|   |       +-- api.py       # Registro de routers
|   |       +-- endpoints/
|   |           +-- auth.py          # /auth: login, refresh
|   |           +-- parquet.py       # /parquet: carga de archivos
|   |           +-- historical.py    # /historical: 1D, 2D, catalogo
|   |           +-- prediction.py    # /prediction: actual e historica
|   |           +-- hydro.py         # /hydro: estaciones IDEAM
|   |           +-- drought.py       # /drought: analisis de sequia
|   |           +-- dashboard.py     # /dashboard: resumen del mapa
|   |           +-- admin_files.py   # /admin/files: CRUD archivos
|   |           +-- admin_datasets.py# /admin/datasets: ciclo datasets
|   |           +-- admin_users.py   # /admin/users: gestion usuarios
|   |           +-- admin_cloud.py   # /admin/cloud: integracion R2
|   |           +-- admin_tiered.py  # /admin/tiered: almac. escalonado
|   |           +-- admin_utils.py   # Utilidades compartidas de admin
|   +-- services/
|   |   +-- historical_data_service.py
|   |   +-- prediction_data_service.py
|   |   +-- hydro_data_service.py
|   |   +-- drought_analysis.py
|   |   +-- ai_summary_service.py
|   |   +-- cache.py
|   |   +-- cloud_storage.py
|   |   +-- parquet_processor.py
|   |   +-- tiered_storage.py
|   +-- schemas/             # Esquemas Pydantic de request/response
+-- .env                     # Variables de entorno (BD, R2, JWT, Redis)
+-- requirements.txt
+-- Dockerfile
+-- docker-compose.yml
\end{verbatim}

\begin{example}
Flujo de una consulta 1D de serie de tiempo:
\begin{enumerate}
    \item El router \texttt{historical.py} recibe \texttt{GET /api/v1/historical/timeseries} con parámetros \texttt{lat}, \texttt{lon}, \texttt{index} y \texttt{scale}.
    \item \texttt{historical\_data\_service.py} verifica el caché (\texttt{cache.py}).
    \item Si no hay caché, DuckDB consulta el archivo Parquet en disco.
    \item El resultado se almacena en caché y se devuelve como JSON.
\end{enumerate}
\end{example}
